%% kapa: exam=1 round=26 year=2015 period=1 form=A subject=civil_law subject_ko=민법 seq=13 number=13
%% source: Q-Net 감정평가사 기출문제 (cst00309) artlSeq=5138001 file=extracted/1008/5138001/제26회_감정평가사 1차_기출/제26회_감정평가사 1차_기출/1교시_감평_A형.pdf
%% conversion: kordoc 4.24.51 -> tools/gen1.py (auto); review=proofread-sonnet; flags=none
\begin{question}{13}
법률행위에 관한 설명으로 옳지 않은 것은? (다툼이 있으면 판례에 따름)
\begin{choices}
  \item 어떠한 의사표시가 진의 아닌 의사표시로서 무효라고 주장하는 자는 그에 대한 증명책임을 진다.
  \item 불공정한 법률행위에 해당하는지는 법률행위가 이루어진 시점을 기준으로 약속된 급부와 반대급부 사이의 객관적 가치를 비교 평가하여 판단해야 한다.
  \item 대리인에 의한 법률행위가 불공정한 법률행위에 해당하는지 판단함에 있어서 궁박은 대리인을 기준으로 한다.
  \item 부동산 매매계약에 있어 쌍방 당사자가 모두 X토지를 계약의 목적물로 삼았으나 그 목적물의 지번 등에 관하여 둘 다 착오를 일으켜 계약서에는 Y토지로 표시한 경우, 매수인은 X토지에 대해 이전등기를 청구할 수 있다.
  \item 특별한 사정이 없는 한, 착오를 이유로 법률행위를 취소하려면 표의자에게 중대한 과실이 없어야 한다.
\end{choices}
\end{question}
\answer{3}
